% Einheit 2 von 4 – Winkel berechnen (bank/trigonometrie/e2.jsonl, zusammenbau v0.1)
%% TODO Zweigzeile Teil 1: was hier gelernt wird (Satz in Schülersprache aus dem Einheitstitel) – Einheit 2
\einheitenkopf[e2]{Einheit 2 von 4 · Winkel berechnen}
\zweigzeile{ab Kl. 10 (am Gymnasium ab Kl. 9) · P10 oft}
\setcounter{aufgabe}{18}

%% TODO Titel: Ich-kann-Satz fehlt in der Bank (Platzhalter: Kettenname) – Nr. 19
%% TODO Anweisung: ein Satz über der ersten Teilaufgabe fehlt in der Bank – Nr. 19
\begin{aufgabe}{Winkel berechnen}
%% TODO \rechenplatz{n}: Schrittzahl des Grundfalls fehlt in der Bank (nur bei fester Schreibform, 2.3 b)
\begin{teile}
\teil Gegeben: Hypotenuse $7$ cm und Gegenkathete $4$ cm. Gesucht: der Winkel $\alpha$. Was ist gesucht? \\ \kreuz{Seite: Gleichung umstellen} \\ \kreuz{Winkel: Umkehrtaste}
\teil Gegenkathete $3$ cm, Hypotenuse $8$ cm – wie groß ist $\alpha$? $\alpha \approx$ \leerfeld $^\circ$
\teil Ankathete $5$ cm, Hypotenuse $7$ cm – wie groß ist $\alpha$? $\alpha \approx$ \leerfeld $^\circ$
\teil Gegenkathete $5$ cm, Ankathete $8$ cm – wie groß ist $\alpha$? $\alpha \approx$ \leerfeld $^\circ$
\teil Rechter Winkel bei $A$, $a = 10$ cm, $b = 3$ cm. Wie groß ist $\beta$?

\dreieck{(0,0)}{(4,0)}{(0,3)}{a}{b}{c}{}{\beta}{}
$\beta \approx$ \leerfeld $^\circ$
\teil Gegenkathete von $\alpha$: $3$ cm, Hypotenuse $5$ cm. Berechne $\alpha$, dann $\beta$ als Ergänzung; kontrolliere $\beta$ mit dem Kosinus. $\alpha \approx$ \leerfeld $^\circ$, $\beta \approx$ \leerfeld $^\circ$
\teil Gegenkathete $2{,}3$ m, Hypotenuse $4{,}7$ m – wie groß ist $\alpha$, auf eine Dezimale? $\alpha \approx$ \leerfeld $^\circ$
\teil Eine Seilbahn überwindet $420$ m Höhe, das Seil ist $1\,300$ m lang. Wie groß ist der Steigungswinkel? \leerfeld $^\circ$
\teil Die Punkte $A$, $B$, $C$ bilden ein Dreieck mit rechtem Winkel bei $A$. Lies die Längen der Katheten ab und berechne den Winkel $\beta$ bei $B$.

\begin{ksys}[xmin=-1,xmax=6,ymin=-1,ymax=5,ablesen] \punkt{1}{1}{A} \punkt{5}{1}{B} \punkt{1}{4}{C} \end{ksys}
$\beta \approx$ \leerfeld $^\circ$
\teil Im Dreieck $ABC$ ist die Höhe von $C$ auf $AB$ $4{,}5$ cm lang, $AC = 7{,}2$ cm. Wie groß ist $\alpha$ bei $A$? $\alpha \approx$ \leerfeld $^\circ$
\teil Ankathete $5{,}2$ m, Hypotenuse $8$ m. Zeige rechnerisch, dass $\alpha$ rund $49^\circ$ groß ist.
\teil Eine Seilbahn führt von der Talstation $A$ zur Bergstation $B$. Das Seil ist $610$ m lang, die waagerechte Strecke unter dem Seil $460$ m. Berechne den Winkel zwischen Seil und Waagerechter. \hfill (P10 2024 OS) \\ \leerfeld $^\circ$
\end{teile}
\end{aufgabe}

%% TODO Titel: Ich-kann-Satz fehlt in der Bank (Platzhalter: Kettenname) – Nr. 20
%% TODO Anweisung: ein Satz über der ersten Teilaufgabe fehlt in der Bank – Nr. 20
\begin{aufgabe}{Winkel prüfen: die längere Kathete liegt dem größeren Winkel gegenüber}
\begin{teile}
\teil Die Katheten sind $4$ cm und $9$ cm lang. Tim sagt: Der Winkel gegenüber der $4$-cm-Kathete ist $66^\circ$. Kann das stimmen?
\end{teile}
\end{aufgabe}

%% TODO Titel: Ich-kann-Satz fehlt in der Bank (Platzhalter: Kettenname) – Nr. 21
%% TODO Anweisung: ein Satz über der ersten Teilaufgabe fehlt in der Bank – Nr. 21
\begin{aufgabe}{Umkehrtaste am Taschenrechner (SHIFT sin, cos, tan; Anzeige sin⁻¹)}
\begin{teile}
\teil Mit der Umkehrtaste: $\mathrm{sin}^{-1}(0{,}42)$ – welcher Winkel, auf eine Dezimale? \leerfeld $^\circ$
\end{teile}
\end{aufgabe}

%% TODO Titel: Ich-kann-Satz fehlt in der Bank (Platzhalter: Kettenname) – Nr. 22
%% TODO Anweisung: ein Satz über der ersten Teilaufgabe fehlt in der Bank – Nr. 22
\begin{aufgabe}{Winkel berechnen – Fehler finden, Begründen, Anwendung}
\begin{teile}
\teil Gegenkathete von $\alpha$: $3$ cm, Ankathete: $8$ cm. Wo ist der Fehler? \rechnung{\mathrm{tan}\,\alpha &= \frac{8}{3} \\ \alpha &\approx 69{,}4^\circ}
\teil Ole tippt $\mathrm{sin}^{-1}(1{,}3)$ und der Taschenrechner meldet einen Fehler. Warum?
\teil Eine Rollstuhlrampe darf höchstens $3{,}4^\circ$ steil sein. Eine $5$ m lange Rampe überwindet $0{,}32$ m. Ist sie zulässig?
\end{teile}
\end{aufgabe}
